\documentclass[12pt]{article}
\usepackage{amssymb,amsmath,pgf,setspace,comment,verbatim,titling,pdflscape,url}
\usepackage{graphicx}
\graphicspath{{../}}
\usepackage[left=1in,right=1in,top=1in,bottom=0.75in]{geometry}
\usepackage[round]{natbib}
\usepackage{hyperref}
\usepackage[labelfont=bf]{caption}
\usepackage{booktabs}
\newcommand{\pderiv}[2]{\frac{\partial#1}{\partial#2}}
\newcommand{\monthyear}{\ifcase\month\or January\or February\or March\or April\or May\or June\or July\or August\or September\or October\or November\or December\fi\ \the\year}
\hypersetup{colorlinks,%
	citecolor=black,%
	filecolor=black,%
	linkcolor=black,%
	urlcolor=blue,%
}
\setstretch{1.5}

\setlength{\droptitle}{-50pt}

\begin{document}
\title{Decision Assistance and Steering in Health Insurance \\
        SUPPLEMENTAL APPENDIX}
\author{Ian McCarthy \& Evan Saltzman}
\date{\monthyear}
\maketitle

\renewcommand{\thetable}{A\arabic{table}}
\renewcommand{\thefigure}{A\arabic{figure}}

\section{Constructing the Outside Option}
\label{app:outside_option}

We construct the outside option from six years (2014--2019) of longitudinal Covered California data, treating a household as uninsured in the years in which it has no exchange coverage but remains eligible for the exchange market. For example, a household with exchange coverage in 2015 and 2016 is treated as uninsured, the outside option, in 2014, 2017, and 2018, provided it remained eligible for the exchange market in those years. A household may lose exchange-market eligibility if it gains access to employer-sponsored insurance or to public coverage such as Medicaid or Medicare. Our data do not record when a household loses eligibility, so we impute it from the Survey of Income and Program Participation (SIPP).

To impute consumer exchange market eligibility, we use data from the SIPP. The SIPP is well-suited for the imputation because (1) it asks the insurance status of respondents for every month over a three-year period (2013-2015) and (2) it includes detailed information on the chief reasons for consumers' coverage status, such as whether the respondent obtained or lost an offer of employer-sponsored insurance, moved in or out of California, or became eligible or ineligible for Medicare or Medicaid. For SIPP respondents who newly obtained or gave up individual market coverage, we construct a \textit{transitioned} variable that indicates whether the respondent gained or lost eligibility for the individual market. The transitioned variable takes value 1 if the respondent (1) belongs to a household that obtained or lost an offer of employer-sponsored insurance; (2) moved out of or into California; (3) turned 65 and qualified for Medicare; and (4) became eligible for Medicaid following a drop in income. We estimate a logit regression of the \textit{transitioned} variable on the demographic variables available in both the SIPP and the Covered California data, including age, income, gender, race, and household size. We use the estimated logit to predict whether Covered California consumers observed for only some years of our study timeframe transitioned into or out of the exchange market. Consumers transitioning into or out of the exchange market are removed from our study population during years when they are not enrolled in an exchange plan.

\section{Assistance Channels}
\label{app:channels}

The main text groups formal decision assistance into two categories, agents and navigators. Covered California's enrollment records distinguish a finer set of service channels, which we map into those categories. Table~\ref{tab:app-channels} lists the underlying codes, a short description of each, the category we group it into, and its share of enrollment records over 2014--2019.

\begin{table}[!h]
\centering
\begin{minipage}[h]{6in}
\caption[]{\textbf{Covered California Assistance Channels}\footnote{Service-channel codes recorded in the Covered California enrollment data, 2014--2019, with each code's description, the higher-level category we group it into, and its share of individual enrollment records. Shares are of all individual enrollment records, assisted and unassisted. The channel shares in the main text are computed over household-years and therefore differ slightly.}}
\label{tab:app-channels}
\centerline{\begin{tabular}{lllr}
\hline\hline
Code & Description & Grouped as & Share \\
\hline
CIA & Certified insurance agent & Agent & 43.5\% \\
PBE & Plan-based enroller & Agent & 1.7\% \\
SCR & Service-center representative & Navigator & 12.2\% \\
CEC & Certified enrollment counselor & Navigator & 6.2\% \\
CEW & Certified enrollment worker & Navigator & 0.4\% \\
Unassisted & Unassisted & Unassisted & 36.1\% \\
\hline\hline
\end{tabular}
}
\end{minipage}
\end{table}

The two commissioned or insurer-affiliated channels are the certified insurance agents and the plan-based enrollers, which together form the agent category. Certified insurance agents are independent private agents who hold appointments with a subset of insurers and are paid a commission by the insurer into which they enroll a consumer. They are the large majority of the agent category and the channel through which the commission steering we study operates. Plan-based enrollers are affiliated with a single health plan and enroll consumers into that insurer's products, so they are a small, structurally distinct group whose influence on insurer choice comes from that affiliation rather than from a commission on competing plans.

The navigator category pools the three non-commissioned channels, none of which is paid by any insurer. Certified enrollment counselors and workers are community-based assisters, and service-center representatives are Covered California's own telephone staff. We group these together because they share the absence of a commission incentive, though they differ in setting, and the service-center channel is the largest of the three.

We use the two-category grouping throughout the paper, because the commission mechanism we model operates at the level of whether a household was assisted by a commissioned agent rather than at the level of the specific channel. The one channel whose distinction bears on our interpretation is the plan-based enroller, since its influence on insurer choice reflects a single-insurer affiliation rather than a commission on competing plans. It is a small share of enrollments and is grouped with the certified agents in the assistance indicator used in estimation.

\subsection{Excluding Plan-Based Enrollers}
\label{app:pbe}

Plan-based enrollers are the one assistance channel that can place a consumer with only a single insurer, so their influence on insurer choice works through a menu restriction, not a commission response. We assess their effect on the demand estimates directly. After merging the raw service-channel code onto the estimation sample by household and year, we find that plan-based enrollers account for 1.3 percent of the household-years in the 20 percent structural sample. Table~\ref{tab:app-pbe} re-estimates the main-text demand specification with these households excluded, alongside the full-sample estimates. Dropping them leaves the price, actuarial value, channel, and nesting parameters essentially unchanged, with no change above 2 percent. The commission-by-agent coefficients, which carry the steering mechanism, move from 0.027 and $-0.066$ in the full sample to 0.030 and $-0.070$ without plan-based enrollers, so the steering is slightly stronger once they are excluded. The single-insurer channel does not drive the estimated steering.

\begin{table}[!h]
\centering
\begin{minipage}[h]{6in}
\caption[]{\textbf{Demand Estimates Excluding Plan-Based Enrollers}\footnote{Main-text demand specification re-estimated on the 20\% structural sample with plan-based-enroller households excluded, alongside the full-sample estimates. Plan-based enrollers are identified by merging the raw Covered California service-channel code onto the estimation households by household and year, and account for 1.3\% of household-years in the sample. Premiums are net of subsidy and measured in hundreds of dollars.}}
\label{tab:app-pbe}
\centerline{\begin{tabular}{lcc}
\hline\hline
 & Full sample & Excl.\ PBE \\
\hline
Premium & -0.415 & -0.416 \\
Actuarial value (AV) & 4.392 & 4.406 \\
Navigator $\times$ AV & 1.185 & 1.181 \\
Agent $\times$ AV & 0.729 & 0.723 \\
Navigator $\times$ premium & -0.098 & -0.097 \\
Agent $\times$ premium & -0.021 & -0.022 \\
Commission $\times$ agent & 0.027 & 0.030 \\
Commission$^2$/100 $\times$ agent & -0.066 & -0.070 \\
$\lambda$ (nesting parameter) & 0.374 & 0.375 \\
\hline\hline
\end{tabular}
}
\end{minipage}
\end{table}

\subsection{Commission Utility and Menu Restriction}
\label{app:menu-sim}

The certified insurance agents present the full menu but may still steer a household within it, short of the plan-based enroller's single-insurer restriction. We represent that steering through the commission coefficient in utility. To check that this does not distort the plan choices the model predicts, we simulate a gatekeeper process and confirm that our specification reproduces its choices.

The simulation draws a market of six insurers whose commissions vary, and generates choices from a process with no commission term in utility. Unassisted households choose their preferred plan from the full menu. Agent-assisted households instead face a restricted menu, in which each insurer appears with a probability that rises in its commission, and choose their preferred plan from what is presented. This produces the steering the commissions are meant to induce, with agent-assisted households moving away from the low-commission insurers and toward the high-commission ones, through availability alone and not through preference.

We then fit our commission-in-utility model to these choices, with the full menu available to every household and the commission entering utility for agent-assisted households. Table~\ref{tab:app-menu-sim} reports the resulting shares. The fitted model recovers a positive commission coefficient and reproduces the gatekeeper agent shares to within 1.1 percentage points, with a correlation of 0.999. A menu restriction and a commission in utility are observationally equivalent for the plan choices the model predicts, and our specification loads the restriction onto the commission coefficient without distorting those predictions.

\begin{table}[!h]
\centering
\begin{minipage}[h]{6in}
\caption[]{\textbf{Menu Restriction and the Commission Coefficient}\footnote{Simulated plan-choice shares. Choices are generated by a gatekeeper process in which agent-assisted households face a menu that includes each insurer with probability rising in its commission, with no commission term in utility. The commission-in-utility model is then fit on the full menu. Columns report the unassisted shares, the gatekeeper agent shares, and the agent shares predicted by the fitted commission model. The fitted commission coefficient is positive, and the fitted agent shares match the gatekeeper shares to within 1.1 percentage points, with a correlation of 0.999.}}
\label{tab:app-menu-sim}
\centerline{\begin{tabular}{rrrrr}
\hline\hline
Insurer & Commission & Unassisted & \shortstack{Broker\\(gatekeeper)} & \shortstack{Broker\\(fitted)} \\
\hline
1 & 0.5 & 0.061 & 0.022 & 0.025 \\
2 & 1.0 & 0.058 & 0.031 & 0.030 \\
3 & 1.5 & 0.010 & 0.008 & 0.007 \\
4 & 2.5 & 0.267 & 0.235 & 0.225 \\
5 & 3.5 & 0.473 & 0.500 & 0.511 \\
6 & 4.5 & 0.131 & 0.205 & 0.202 \\
\hline\hline
\end{tabular}
}
\end{minipage}
\end{table}

The baseline simulation lets the menu depend only on commissions. In practice, which plans an agent presents may be correlated with the agent or with the households the agent serves, so the menu restriction could be selected on characteristics that also shape plan choice. We extend the simulation to allow this. Menu availability now depends on an unobserved household type that also shifts utility, capturing either the agent a household is matched to or the way households sort across agents, and we vary the strength of that selection. Table~\ref{tab:app-gatekeeper-selection} reports the fit as the selection strengthens, both when the selection is aligned with commissions, the realistic case in which agents who restrict menus also favor high-commission insurers, and when it is orthogonal to commissions, an adversarial case. Across the range the commission-in-utility model continues to reproduce the gatekeeper agent shares to within 1.7 percentage points, with correlations above 0.997. The equivalence is not an artifact of random menus. It survives selection of the menu on agent and household characteristics, because the plan fixed effects absorb the average tilt while the commission coefficient absorbs the part that covaries with commissions.

\begin{table}[!h]
\centering
\begin{minipage}[h]{6in}
\caption[]{\textbf{Menu Restriction with Agent and Household Selection}\footnote{Simulated plan choices under a gatekeeper process in which the menu is selected on an unobserved household type that also shifts utility. Selection strength zero is the random-gatekeeper baseline of Table~\ref{tab:app-menu-sim}. Aligned selection targets high-commission insurers; orthogonal selection targets a dimension uncorrelated with commissions. Each row fits the commission-in-utility model and reports the fitted commission coefficient, the maximum absolute difference between the gatekeeper and fitted agent shares, and their correlation.}}
\label{tab:app-gatekeeper-selection}
\centerline{\begin{tabular}{lrrrr}
\hline\hline
Selection & Strength & Commission coef. & Max abs.\ diff.\ & Correlation \\
\hline
Aligned with commission & 0.0 & 0.306 & 0.0100 & 0.9996 \\
 & 0.5 & 0.252 & 0.0122 & 0.9996 \\
 & 1.0 & 0.153 & 0.0087 & 0.9996 \\
 & 2.0 & 0.035 & 0.0171 & 0.9974 \\
\hline
Orthogonal to commission & 0.0 & 0.296 & 0.0106 & 0.9996 \\
 & 0.5 & 0.281 & 0.0122 & 0.9995 \\
 & 1.0 & 0.247 & 0.0119 & 0.9994 \\
 & 2.0 & 0.187 & 0.0144 & 0.9984 \\
\hline\hline
\end{tabular}
}
\end{minipage}
\end{table}

The equivalence concerns the positive predictions and the demand fit. It does not carry over to the counterfactuals. A menu restriction fixed by an agent's appointments does not respond to a change in commissions, while a commission term in utility does, so the two models diverge once we alter commissions. In the counterfactuals we therefore attribute the steering to the commission margin, as the main text describes. For the demand estimates, the exercise shows that representing the restriction through the commission coefficient does not distort the plan choices the model predicts.

\section{Additional Design-Based Results}
\label{app:reduced_form}

The design-based estimates in the main text come from a prediction-based estimator. We fit the outcome model on unassisted households, weighted by inverse propensity weights, predict the counterfactual outcome for assisted households, and form the ATT as the difference between what assisted households chose and what the model predicts they would have chosen. Assistance never appears as a regressor, so there is no assistance coefficient to report. This appendix reports the complementary pooled specifications, in which assistance enters as a covariate, so its coefficient can be reported across a sequence of control sets.

\subsection{Dominated Choices}

Table~\ref{tab:app-dominated} reports linear probability models of the dominated-choice indicator on the sample of households with at least one dominated plan in the choice set. Column (1) is unconditional. Column (2) adds the household demographics that enter the propensity score, namely age composition, gender, race and ethnicity, and household size. Income does not appear, because the sample is defined by cost-sharing-reduction eligibility and therefore contains only households below 200\% of the federal poverty level, leaving the income brackets used elsewhere in the paper with no variation to exploit. Column (3) adds a new-enrollee indicator and year fixed effects, and column (4) adds region fixed effects. Standard errors are clustered by rating region throughout, and all specifications are inverse-propensity weighted.

The assistance coefficient is stable across the columns, moving only from $-0.0152$ to $-0.0168$ as controls accumulate and settling at $-0.0161$ once region fixed effects absorb differences in market composition.

Table~\ref{tab:app-dominated-sample} reports the main-text estimator re-run on new enrollees alongside the all-enrollee estimates. The main-text sample is all enrollees so that the design-based analysis and the structural model are estimated on the same households. Restricting to new enrollees produces larger effects for every channel, consistent with assistance mattering most for households making a first plan choice, and leaves the ordering across channels unchanged.

\begin{table}[!h]
\centering
\begin{minipage}[h]{6in}
\caption[]{\textbf{Dominated Choices, Pooled Specifications}\footnote{Linear probability models of an indicator for selecting a dominated plan, on the sample of households with at least one dominated plan in the choice set. All specifications are weighted by inverse propensity weights, with standard errors clustered by rating region. Demographics include age composition, gender, race and ethnicity, and household size. $^{***}p<0.01$, $^{**}p<0.05$, $^{*}p<0.10$.}}
\label{tab:app-dominated}
\resizebox{\linewidth}{!}{%
    
\begin{tabular}{lcccc}
\toprule
Variable & (1) & (2) & (3) & (4)\\
\midrule
Assisted & -0.0152*** & -0.0152*** & -0.0168*** & -0.0161***\\
 & (0.0012) & (0.0011) & (0.0011) & (0.0007)\\
HH Size &  & -0.0018*** & -0.0007 & -0.0003\\
 &  & (0.0004) & (0.0005) & (0.0005)\\
New Enrollee &  &  & 0.0168*** & 0.0168***\\
 &  &  & (0.0027) & (0.0027)\\
Demographics &  & X & X & X\\
Year FE &  &  & X & X\\
Region FE &  &  &  & X\\
Observations & 2,993,294 & 2,993,294 & 2,993,294 & 2,993,294\\
R$^2$ & 0.002 & 0.003 & 0.012 & 0.016\\
\bottomrule
\end{tabular}

}
\end{minipage}
\end{table}

\begin{table}[!h]
\centering
\begin{minipage}[h]{6in}
\caption[]{\textbf{Dominated Choice ATT, All Enrollees Versus New Enrollees}\footnote{Prediction-based ATT of decision assistance on the probability of selecting a dominated plan, weighted by inverse propensity weights. The all-enrollee column is the main-text estimate. The new-enrollee column re-runs the same estimator on households enrolling for the first time.}}
\label{tab:app-dominated-sample}
\centerline{%
    
\begin{tabular}{lcc}
\toprule
Channel & All Enrollees & New Enrollees\\
\midrule
Any Assistance & -0.0156 & -0.0226\\
Agent/Broker & -0.0183 & -0.0252\\
Navigator & -0.0090 & -0.0180\\
\bottomrule
\end{tabular}

}
\end{minipage}
\end{table}

\clearpage

\subsection{Plan Choice}

Table~\ref{tab:app-choice} reports the companion sequence for plan choice. The model is a conditional logit over the plans in each household's choice set, pooled across assisted and unassisted households and weighted by inverse propensity weights. Assistance enters through its interaction with metal tier, separately by channel, using the same navigator and agent split as the structural model in the main text. Column (1) includes plan attributes only. Column (2) adds the premium-by-demographic interactions that constitute the full control set in the main-text specification.

Two features of this model differ from the dominated-choice table. First, there is no region-fixed-effects column. A household-level region indicator is constant across the alternatives in that household's choice set and is differenced out of the conditional logit, so it does not control for anything. The term that would actually absorb region here is region interacted with metal tier and insurer, which adds well over a hundred parameters and leaves the insurer block empty wherever an insurer does not operate in a region. We do not estimate it, so region is uncontrolled in this model; it is not absorbed.

Second, assistance interacts only with silver and bronze. Catastrophic plans are excluded from the sample, so the four metal tiers exhaust the alternatives and a full set of four assistance-by-metal interactions sums to a household constant, which is differenced out. Gold and platinum share the omitted category, matching the base specification.

The steering pattern, toward silver with agents steering more strongly than navigators, is stable across both columns.

\begin{table}[!h]
\centering
\begin{minipage}[h]{6in}
\caption[]{\textbf{Plan Choice, Pooled Specifications}\footnote{Conditional logit over the plans in each household's choice set, pooled across assisted and unassisted households and weighted by inverse propensity weights. Premiums are net of subsidy and measured in hundreds of dollars. Plan attributes are metal tier, network type, health savings account eligibility, and insurer indicators. Standard errors in parentheses.}}
\label{tab:app-choice}
\centerline{%
    
\begin{tabular}{lcc}
\toprule
Variable & (1) & (2)\\
\midrule
Premium & -1.2004 (0.0023) & -0.7460 (0.0045)\\
Navigator $\times$ Silver & 0.2942 (0.0094) & 0.2821 (0.0095)\\
Navigator $\times$ Bronze & -0.2836 (0.0105) & -0.3113 (0.0106)\\
Agent $\times$ Silver & 0.6093 (0.0077) & 0.6157 (0.0078)\\
Agent $\times$ Bronze & 0.0721 (0.0084) & 0.0909 (0.0086)\\
Premium $\times$ demographics &  & X\\
Households & 1,067,827 & 1,067,827\\
Log-likelihood & -2,603,785 & -2,544,116\\
\bottomrule
\end{tabular}

}
\end{minipage}
\end{table}

\section{Model Formulas}\label{app:formulas}

This appendix summarizes the formulas used in our model.   Household $i$ in market $m$ and year $t$ chooses among the exchange plans $j \in J_{mt}$ and the outside option, with deterministic utilities $V_{ijt}$ and $V_{i0t} = \alpha_{it}\, \rho_{it}$, nesting parameter $\lambda$, and household price coefficient $\alpha_{it}$.

\paragraph{Choice probabilities.} The choice probabilities in the main text are

\begin{equation*}
	q_{ijt}({\bf p},{\bf y}) =  \frac{e^{V_{ijt}({\bf p},{\bf y})/\lambda}\left(\sum_{k \in J_{mt}} e^{V_{ikt}({\bf p},{\bf y})/\lambda}\right)^{\lambda-1}}{e^{V_{i0t}} + \left(\sum_{k \in J_{mt}} e^{V_{ikt}({\bf p},{\bf y})/\lambda}\right)^{\lambda}},
\end{equation*}


\noindent Let $I_{it}({\bf p},{\bf y}) \equiv \log \sum_{k \in J_{mt}} e^{V_{ikt}({\bf p},{\bf y})/\lambda}$ denote the inclusive value of the exchange nest. The within-nest and nest probabilities are
\begin{equation}
\label{eqn:app-shares}
q'_{ijt}({\bf p},{\bf y}) = \frac{e^{V_{ijt}({\bf p},{\bf y})/\lambda}}{\sum_{k \in J_{mt}} e^{V_{ikt}({\bf p},{\bf y})/\lambda}}, \qquad
q^{\,\mathrm{nest}}_{it}({\bf p},{\bf y}) = \frac{e^{\lambda I_{it}({\bf p},{\bf y})}}{e^{\lambda I_{it}({\bf p},{\bf y})} + e^{V_{i0t}}},
\end{equation}
and the plan choice probability is $q_{ijt}({\bf p},{\bf y}) = q'_{ijt}({\bf p},{\bf y})\, q^{\,\mathrm{nest}}_{it}({\bf p},{\bf y})$, which is the nested-logit expression in the main text. Total plan demand is the sum of the choice probabilities, i.e.,  $q_{jt}({\bf p},{\bf y}) = \sum_{i \in I} \, q_{ijt}({\bf p},{\bf y})$.

\paragraph{Price and commission derivatives.} Under the two-part form in the main text, the within-nest probability $q'_{ijt}$ is evaluated at the household's realized channel, while the nest probability uses the expected inclusive value $\bar{I}_{it} = \sum_c \pi_{ict} I^c_{it}$, with $I^c_{it}$ the inclusive value under channel state $c$ and $q'^{\,c}_{ijt}$ the within-nest probability under that state. The partial derivative of demand with respect to the insurer's base premium is

\begin{equation*}
		\frac{\partial q_{ikt}({\bf p},{\bf y})}{\partial p_{jmt}} = \sum_{l \in J_{mt}} \frac{ \partial q_{ikt}({\bf p},{\bf y})}{\partial p_{ilt}({\bf p})}\frac{ \partial p_{ilt}({\bf p})}{\partial p_{jmt}}
	\end{equation*}

	\noindent for all plans $j,k$, where $\partial p_{ilt}/\partial p_{jmt}$ carries the household's age rating factor and the subsidy rule in the main text, and

\begin{equation*}
		\frac{\partial q_{ikt}({\bf p},{\bf y})}{\partial  p_{ijt}} = q_{ikt}({\bf p},{\bf y}) \left[ \frac{\alpha_{it}}{\lambda}\left( \mathbb{I}[k = j] - q'_{ijt}({\bf p},{\bf y}) \right) + \left(1 - q^{\,\mathrm{nest}}_{it}({\bf p},{\bf y})\right) \sum_c \pi_{ict}\, \alpha^c_{it}\, q'^{\,c}_{ijt}({\bf p},{\bf y}) \right],
\end{equation*}

\noindent with $\alpha_{it}$ the premium coefficient at the household's realized channel and $\alpha^c_{it}$ the coefficient under channel state $c$. The first term moves the household among plans and the second moves it into or out of the market through the expected inclusive value. The partial derivative of demand with respect to the commission is

\begin{equation*}
	\begin{split}
		\frac{\partial q_{ikt}({\bf p},{\bf y})}{\partial  y_{jt}} = {} & q_{ikt}({\bf p},{\bf y})\, \mu_{ijt}\, \phi_{ijt} \left[ \frac{a^B_{it}}{\lambda}\left( \mathbb{I}[k = j] - q'_{ijt}({\bf p},{\bf y}) \right) \right. \\
		& \left. {} + \left(1 - q^{\,\mathrm{nest}}_{it}({\bf p},{\bf y})\right) \pi_{iAt}\, q'^{\,A}_{ijt}({\bf p},{\bf y}) \right], \qquad \mu_{ijt} = \beta^{y}_{1} + 2\beta^{y}_{2}\, \phi_{ijt} y_{jt},
	\end{split}
\end{equation*}

\noindent where $\phi_{ijt}$ is the ratio of the household's commission to the plan's commission, equal to one under a flat schedule and to the household's rating ratio under a percentage schedule, and $A$ denotes the agent state. The factor $\mu_{ijt}$ is the household's marginal utility of a commission dollar, which under the quadratic commission terms depends on the household's own commission $\phi_{ijt} y_{jt}$. The plan-choice term applies only to agent-assisted households. The enrollment term applies to every household, since each household's enrollment depends on the agent state with weight $\pi_{iAt}$.

\paragraph{Profit terms.}  Total premium revenue collected by firm $f$ at time $t$ is given by

\begin{equation*}
		R_{ft}({\bf p},{\bf y}) = \sum_{i \in I, m \in M, k \in J_{fmt}} \mathbb{I}_{i,m,t} \sigma_{it} p_{kmt} q_{ikt}({\bf p},{\bf y})
\end{equation*}

\noindent where $\mathbb{I}_{i,m,t}$ indicates household $i$ lives in market $m$ at time $t$.  Total claims incurred by firm $f$ at time $t$ is given by

\begin{equation*}
		C_{ft}({\bf p},{\bf y}) = \sum_{m \in M, k \in J_{fmt}}  c_{kmt}({\bf p},{\bf y}) q_{kmt}({\bf p},{\bf y})
\end{equation*}

\noindent where $q_{kmt}({\bf p},{\bf y}) = \sum_{i \in I} \mathbb{I}_{i,m,t} q_{ikt}({\bf p},{\bf y})$.  Total commissions paid by firm $f$ at time $t$ is given by

\begin{equation*}
    Y_{ft}({\bf p},{\bf y}) = \sum_{j \in J_{ft}} \sum_{i \in I} a_{it}^B \, y_{jt}  \, q_{ijt}({\bf p},{\bf y}) 
\end{equation*}

\noindent In our notation, the official risk adjustment transfer  $RA_{ft}({\bf p},{\bf y})$ as derived by \citet{pope2014} equals
	
	
	
		\begin{equation*}\label{ra_formula}
			RA_{ft}({\bf p},{\bf y}) = \sum_{m \in M, k \in J_{fmt}} \left[\widehat{c}_{kmt}({\bf p},{\bf y}) - \widetilde{c}_{kmt}({\bf p},{\bf y})\right] q_{kmt}({\bf p},{\bf y}) 
		\end{equation*}
		

	
	
	\noindent where the plan's expected average claims with adverse selection $\widehat{c}_{kmt}({\bf p},{\bf y})$ and the plan's expected average claims without adverse selection $\widetilde{c}_{kmt}({\bf p},{\bf y})$ are given by 
    
        \begin{equation*}
            \widehat{c}_{kmt}({\bf p},{\bf y}) = \frac{\widehat{h}_{kmt}({\bf p},{\bf y})}{\sum_{n \in M,l \in J_{nt}} \widehat{h}_{lnt}({\bf p},{\bf y})q_{lnt}({\bf p},{\bf y})} \, \nu\overline{p}
        \end{equation*}

        \begin{equation*}
            \widetilde{c}_{kmt}({\bf p},{\bf y}) = \frac{\widetilde{h}_{kmt}({\bf p},{\bf y})}{\sum_{n \in M, l \in J_{nt}} \widetilde{h}_{lnt}({\bf p},{\bf y})q_{lnt}({\bf p},{\bf y})} \, \nu\overline{p}
        \end{equation*}
    
    
    \noindent The average statewide premium is $\overline{p}$ and the expected percentage of collected premiums that is spent on claims is $\nu$.\footnote{From 2014-2017, $\nu$ was set to 100\% in the ACA transfer formula.  Starting in 2018, $\nu$ was set to 86\%.}  The cost factor $\widehat{h}_{jmt}({\bf p},{\bf y}) \equiv \mathrm{IDF}_j\mathrm{GCF}_{mt}r_{jmt}({\bf p},{\bf y})$ is the product of the plan's induced demand factor,  geographic cost factor, and risk score. The cost factor $\widetilde{h}_{jmt}({\bf p},{\bf y}) \equiv \mathrm{AV}_j\mathrm{IDF}_j\mathrm{GCF}_{mt}a_{jmt}({\bf p},{\bf y})$ is the product of the plan's AV, induced demand factor, geographic cost factor, and average ACA age rating factor $a_{jmt}({\bf p},{\bf y}) \equiv  \frac{\sum_{i \in I} \left(\mathbb{I}_{i,m,t} \right) a_{it} q_{ijt}({\bf p},{\bf y})}{q_{jmt}({\bf p},{\bf y})}$  across the plan's enrollees, where $a_{it}$ is the household's CMS age rating factor. The induced demand factor accounts for expected moral hazard associated with lower cost sharing when enrolled in a more generous plan. CMS sets $\mathrm{IDF}_j$ equal to 1 for bronze plans, 1.03 for silver plans, 1.08 for gold plans, and 1.15 for platinum plans.  The geographic cost factors account for regional differences in the cost of doing business (e.g., input prices, provider market power, etc.)  CMS publishes geographic cost factors in its risk adjustment reports.  The plan risk score is the fitted value of the risk-score equation in the main text.  The ACA's transitional reinsurance program, active between 2014-2016, helps to offset the cost of an insurer's high-spending enrollees.  Because we lack the individual-level claims data to precisely calculate reinsurance payments, we instead calculate the actuarial value of the reinsurance contract $\iota_{ft}$ using CMS rate filings.  Reinsurance received is then

    	\begin{equation*}\label{reins_formula}
			RI_{ft}({\bf p},{\bf y}) = \iota_{ft} C_{ft}({\bf p},{\bf y})
		\end{equation*}

    \noindent Finally, variable administrative cost is given by

        \begin{equation*}
            VC_{ft}({\bf p},{\bf y}) = vc_{ft} q_{ft}({\bf p},{\bf y})
        \end{equation*}
    
    \noindent where $vc_{ft}$ is per-capita variable administrative cost and $q_{ft}({\bf p},{\bf y}) = \sum_{k \in J_f} q_{kmt}({\bf p},{\bf y})$ is firm demand.
    




\paragraph{Premium first-order conditions.} The partial derivative of firm revenue with respect to the base premium is

    \begin{equation*}
		\frac{\partial R_{ft}({\bf p},{\bf y})}{\partial p_{jmt}}  =   \sum_{i \in I, k \in J_{fmt}} \mathbb{I}_{i,m,t} \sigma_{it} \left(    q_{ijt}({\bf p},{\bf y}) \, \, +   p_{kmt} \frac{ \partial q_{ikt}({\bf p},{\bf y})}{\partial p_{jmt}}\right)
	\end{equation*}

\noindent The partial derivative of firm claims with respect to the base premium is

	\begin{equation*}
		\frac{\partial C_{ft}({\bf p},{\bf y})}{\partial p_{jmt}} =   \sum_{k \in J_{fmt}}  \left[ c_{kmt}({\bf p},{\bf y})\frac{ \partial q_{kmt}({\bf p},{\bf y})}{\partial p_{jmt}}   +  q_{kmt}({\bf p},{\bf y})\frac{ \partial c_{kmt}({\bf p},{\bf y})}{\partial p_{jmt}}  \right] 	
	\end{equation*}
    
\noindent where

\begin{equation*}
	\frac{\partial c_{kmt}({\bf p},{\bf y})}{\partial p_{jmt}} 
	= \mu_t^r \frac{c_{kmt}({\bf p},{\bf y})}{ q_{kmt}({\bf p},{\bf y})}  \sum_{d \in D} \gamma_t^d \left[\frac{\partial q_{dkmt}({\bf p},{\bf y})}{\partial p_{jmt}} - s_{dkmt}({\bf p},{\bf y}) \frac{\partial q_{kmt}({\bf p},{\bf y})}{\partial p_{jmt}} \right] 
\end{equation*}

\noindent The partial derivative of total firm commissions with respect to the base premium is

\begin{equation*}
    \frac{ \partial Y_{ft}({\bf p},{\bf y})}{\partial p_{jmt}} = \sum_{j \in J_{ft}} \sum_{i \in I} a_{it}^B \, y_{jt}  \, \frac{ \partial q_{kmt}({\bf p},{\bf y})}{\partial p_{jmt}}
\end{equation*}

\noindent 	Define $ rs_{kmt}({\bf p},{\bf y}) \equiv \frac{\widehat{h}_{kmt}({\bf p},{\bf y})q_{kmt}({\bf p},{\bf y})}{\sum_{n \in M,l \in J_{nt}} \widehat{h}_{lnt}({\bf p},{\bf y})q_{lnt}({\bf p},{\bf y})}$ and $us_{kmt}({\bf p},{\bf y}) \equiv \frac{\widetilde{h}_{kmt}({\bf p},{\bf y})q_{kmt}({\bf p},{\bf y})}{\sum_{n \in M, l \in J_{nt}} \widetilde{h}_{lnt}({\bf p},{\bf y})q_{lnt}({\bf p},{\bf y})}$.  Let $R_t({\bf p},{\bf y}) \equiv  \sum_{f \in F} R_{ft}({\bf p},{\bf y})$ be total premiums collected. % compute_ra_foc in helpers/ra.R evaluates this derivative in full at the
% carrier level: the share response, the response of each plan's risk score
% and average rating factor to its enrollee mix, and the feedback through
% premiums collected. The premium total R_t is held at premiums collected
% at the observed premiums, as in Evan's code, so it enters the level as a
% constant and the derivative through its own response. The claims
% derivative above carries the matching mix term (compute_claims_comp). The
% analytic conditions are checked against the numerical derivative of the
% model's own profit (scratch/foc_vs_numeric.R, 2026-09-19).
The partial derivative of the risk adjustment transfer with respect to the base premium is

\begin{footnotesize}
		
		\begin{eqnarray*}
			\frac{\partial RA_{ft}({\bf p},{\bf y})}{\partial p_{jmt}} &=& \nu \sum_{k \in J_{fmt}}  \left[\frac{\partial R_t({\bf p},{\bf y})}{\partial p_{jmt}}\left( rs_{kmt}({\bf p},{\bf y}) - us_{kmt}({\bf p},{\bf y})  \right) \right.\\
			&& + \,\,\, \left. R_t({\bf p},{\bf y}) \left( \frac{\partial rs_{kmt}({\bf p},{\bf y})}{\partial p_{jmt}} - \frac{\partial  us_{kmt}({\bf p},{\bf y})}{\partial p_{jmt}} \right) \right]
		\end{eqnarray*}
		
	\end{footnotesize}


\noindent where
	
	
	
	
	
	
	
	\begin{footnotesize}
		\setlength{\arraycolsep}{1.5pt}
		
		\begin{eqnarray*}
			\frac{ \partial rs_{kmt}({\bf p},{\bf y})}{\partial p_{jmt}} &=& \left(\sum_{n\in M,l \in J_{nt}} \widehat{h}_{lnt}({\bf p},{\bf y}) q_{lnt}({\bf p},{\bf y})\right)^{-1}\left[\left(\widehat{h}_{kmt}({\bf p},{\bf y})\frac{\partial q_{kmt}({\bf p},{\bf y})}{\partial p_{jmt}} + q_{kmt}({\bf p},{\bf y})\frac{\partial \widehat{h}_{kmt}({\bf p},{\bf y})}{\partial p_{jmt}} \right) \right.\\
			&& - \,\,\, \left. rs_{kmt}({\bf p},{\bf y}) \sum_{l \in J_{mt}}  \left(\widehat{h}_{lmt}({\bf p},{\bf y})\frac{\partial q_{lmt}({\bf p},{\bf y})}{\partial p_{jmt}} + q_{lmt}({\bf p},{\bf y})\frac{\partial \widehat{h}_{lmt}({\bf p},{\bf y})}{\partial p_{jmt}} \right)\right] 
		\end{eqnarray*}	
		
	\end{footnotesize}
	
	
	\begin{footnotesize}
		\setlength{\arraycolsep}{1.5pt}
		
		\begin{eqnarray*}
			\frac{ \partial us_{kmt}({\bf p},{\bf y})}{\partial p_{jmt}} &=& \left(\sum_{n\in M,l \in J_{nt}} \widetilde{h}_{lnt}({\bf p},{\bf y}) q_{lnt}({\bf p},{\bf y})\right)^{-1}\left[\left(\widetilde{h}_{kmt}({\bf p},{\bf y})\frac{\partial q_{kmt}({\bf p},{\bf y})}{\partial p_{jmt}} + q_{kmt}({\bf p},{\bf y})\frac{\partial \widetilde{h}_{kmt}({\bf p},{\bf y})}{\partial p_{jmt}} \right) \right.\\
			&& - \,\,\, \left. us_{kmt}({\bf p},{\bf y}) \sum_{l \in J_{mt}}  \left(\widetilde{h}_{lmt}({\bf p},{\bf y})\frac{\partial q_{lmt}({\bf p},{\bf y})}{\partial p_{jmt}} + q_{lmt}({\bf p},{\bf y})\frac{\partial\widetilde{h}_{lmt}({\bf p},{\bf y})}{\partial p_{jmt}} \right)\right] 
		\end{eqnarray*}	
		
	\end{footnotesize}

\noindent and

\begin{eqnarray*}
		\frac{\partial \widehat{h}_{kmt}({\bf p},{\bf y})}{\partial p_{jmt}} &=& \mathrm{IDF}_k\mathrm{GCF}_{mt} \left[\frac{r_{kmt}({\bf p},{\bf y})}{q_{kmt}({\bf p},{\bf y})} \sum_{d \in D} \gamma^d \left[\frac{\partial q_{dkmt}({\bf p},{\bf y})}{\partial p_{jmt}} - s_{dkmt}({\bf p},{\bf y}) \frac{\partial q_{kmt}({\bf p},{\bf y})}{\partial p_{jmt}} \right] \right]
	\end{eqnarray*}
	
	\vspace{-0.3in}
	
	\begin{eqnarray*}
		\frac{\partial \widetilde{h}_{kmt}({\bf p},{\bf y})}{\partial p_{jmt}} &=& \mathrm{AV}_k \mathrm{IDF}_k\mathrm{GCF}_{mt} \frac{\sum_{i \in I} \left(\mathbb{I}_{i,m,t} \right) a_{it} \frac{ \partial q_{ikt}({\bf p},{\bf y})}{\partial p_{jmt}} - a_{kmt}({\bf p},{\bf y}) \frac{\partial q_{kmt}({\bf p},{\bf y})}{\partial p_{jmt}}}{q_{kmt}({\bf p},{\bf y})}
	\end{eqnarray*}


\paragraph{Commission first-order condition.} Differentiating the same variable profit with respect to a commission-setting unit's scale $k$ gives the commission first-order condition in the main text, one condition per unit, where a unit is the insurer or, where the filed schedule sets distinct network rates, the insurer-network. Its left-hand side is the marginal benefit of a higher commission, the additional variable profit from the agent-assisted enrollment the commission attracts, accumulated over all of the insurer's plans through the steering coefficients $\beta^{y}_{1}$ and $\beta^{y}_{2}$, together with the induced change in risk-adjustment transfers. Its right-hand side is the marginal outlay on the unit's own schedule, $\sum_{j} w_{jt}\, q^{B}_{jt}$ over the unit's plans, where $w_{jt}$ is the plan's commission basis. The condition sets the ratio of marginal benefit to marginal outlay equal to $(1 - \beta_f) + \delta_f$, where $\beta_f$ is the administrative saving per commission dollar and $\delta_f$ is a carrier-level constant, one level per insurer shared across its units and estimated with the cost parameters. The counterfactuals solve the same condition for each unit's scale.

The partial derivative of firm revenue with respect to the commission is

\begin{equation*}
		\frac{\partial R_{ft}({\bf p},{\bf y})}{\partial y_{jt}} = \sum_{i \in I, m \in M, k \in J_{fmt}} \mathbb{I}_{i,m,t} \sigma_{it} p_{kmt} \frac{\partial q_{ikt}({\bf p},{\bf y})}{\partial y_{jt}}
\end{equation*}

\noindent The partial derivative of firm claims with respect to the commission is

\begin{equation*}
		\frac{\partial C_{ft}({\bf p},{\bf y})}{\partial y_{jt}} = \sum_{m \in M, k \in J_{fmt}}  \left[\frac{\partial c_{kmt}({\bf p},{\bf y})}{\partial y_{jt}} q_{kmt}({\bf p},{\bf y}) + c_{kmt}({\bf p},{\bf y}) \frac{\partial q_{ikt}({\bf p},{\bf y})}{\partial y_{jt}}\right]
\end{equation*}

\noindent where

\begin{equation*}
	\frac{\partial c_{kmt}({\bf p},{\bf y})}{\partial y_{jt}} 
	= \mu_t^r \frac{c_{kmt}({\bf p},{\bf y})}{ q_{kmt}({\bf p},{\bf y})}  \sum_{d \in D} \gamma_t^d \left[\frac{\partial q_{dkmt}({\bf p},{\bf y})}{\partial y_{jt}} - s_{dkmt}({\bf p},{\bf y}) \frac{\partial q_{kmt}({\bf p},{\bf y})}{\partial y_{jt}} \right] 
\end{equation*}

\noindent The partial derivative of total firm commissions paid with respect to the commission is

\begin{equation*}
    \frac{\partial Y_{ft}({\bf p},{\bf y})}{\partial y_{jt}} = \sum_{j \in J_{ft}} \sum_{i \in I} a_{it}^B\left[  \, q_{ijt}({\bf p},{\bf y}) +  y_{jt}  \, \frac{\partial q_{kmt}({\bf p},{\bf y})}{\partial y_{jt}} \right]
\end{equation*}

\noindent The partial derivative of the risk adjustment transfer with respect to the commission is

\begin{footnotesize}
		
		\begin{eqnarray*}
			\frac{\partial RA_{ft}({\bf p},{\bf y})}{\partial y_{jt}} &=& \nu \sum_{k \in J_{fmt}}  \left[\frac{\partial R_t({\bf p},{\bf y})}{\partial  y_{jt}}\left( rs_{kmt}({\bf p},{\bf y}) - us_{kmt}({\bf p},{\bf y})  \right) \right.\\
			&& + \,\,\, \left. R_t({\bf p},{\bf y}) \left( \frac{\partial rs_{kmt}({\bf p},{\bf y})}{\partial  y_{jt}} - \frac{\partial  us_{kmt}({\bf p},{\bf y})}{\partial  y_{jt}} \right) \right]
		\end{eqnarray*}
		
	\end{footnotesize}

\noindent where

\begin{footnotesize}
		\setlength{\arraycolsep}{1.5pt}
		
		\begin{eqnarray*}
			\frac{ \partial rs_{kmt}({\bf p},{\bf y})}{\partial y_{jt}} &=& \left(\sum_{n\in M,l \in J_{nt}} \widehat{h}_{lnt}({\bf p},{\bf y}) q_{lnt}({\bf p},{\bf y})\right)^{-1}\left[\left(\widehat{h}_{kmt}({\bf p},{\bf y})\frac{\partial q_{kmt}({\bf p},{\bf y})}{\partial y_{jt}} + q_{kmt}({\bf p},{\bf y})\frac{\partial \widehat{h}_{kmt}({\bf p},{\bf y})}{\partial y_{jt}} \right) \right.\\
			&& - \,\,\, \left. rs_{kmt}({\bf p},{\bf y}) \sum_{l \in J_{mt}}  \left(\widehat{h}_{lmt}({\bf p},{\bf y})\frac{\partial q_{lmt}({\bf p},{\bf y})}{\partial y_{jt}} + q_{lmt}({\bf p},{\bf y})\frac{\partial \widehat{h}_{lmt}({\bf p},{\bf y})}{\partial y_{jt}} \right)\right] 
		\end{eqnarray*}	
		
	\end{footnotesize}
	
	
	\begin{footnotesize}
		\setlength{\arraycolsep}{1.5pt}
		
		\begin{eqnarray*}
			\frac{ \partial us_{kmt}({\bf p},{\bf y})}{\partial y_{jt}} &=& \left(\sum_{n\in M,l \in J_{nt}} \widetilde{h}_{lnt}({\bf p},{\bf y}) q_{lnt}({\bf p},{\bf y})\right)^{-1}\left[\left(\widetilde{h}_{kmt}({\bf p},{\bf y})\frac{\partial q_{kmt}({\bf p},{\bf y})}{\partial y_{jt}} + q_{kmt}({\bf p},{\bf y})\frac{\partial \widetilde{h}_{kmt}({\bf p},{\bf y})}{\partial y_{jt}} \right) \right.\\
			&& - \,\,\, \left. us_{kmt}({\bf p},{\bf y}) \sum_{l \in J_{mt}}  \left(\widetilde{h}_{lmt}({\bf p},{\bf y})\frac{\partial q_{lmt}({\bf p},{\bf y})}{\partial y_{jt}} + q_{lmt}({\bf p},{\bf y})\frac{\partial\widetilde{h}_{lmt}({\bf p},{\bf y})}{\partial y_{jt}} \right)\right] 
		\end{eqnarray*}	
		
	\end{footnotesize}

\noindent and

\begin{eqnarray*}
		\frac{\partial \widehat{h}_{kmt}({\bf p},{\bf y})}{\partial y_{jt}} &=& \mathrm{IDF}_k\mathrm{GCF}_{mt} \left[\frac{r_{kmt}({\bf p},{\bf y})}{q_{kmt}({\bf p},{\bf y})} \sum_{d \in D} \gamma^d \left[\frac{\partial q_{dkmt}({\bf p},{\bf y})}{\partial y_{jt}} - s_{dkmt}({\bf p},{\bf y}) \frac{\partial q_{kmt}({\bf p},{\bf y})}{\partial y_{jt}} \right] \right]
	\end{eqnarray*}
	
	\vspace{-0.3in}
	
	\begin{eqnarray*}
		\frac{\partial \widetilde{h}_{kmt}({\bf p},{\bf y})}{\partial y_{jt}} &=& \mathrm{AV}_k \mathrm{IDF}_k\mathrm{GCF}_{mt} \frac{\sum_{i \in I} \left(\mathbb{I}_{i,m,t} \right) a_{it} \frac{ \partial q_{ikt}({\bf p},{\bf y})}{\partial y_{jt}} - a_{kmt}({\bf p},{\bf y}) \frac{\partial q_{kmt}({\bf p},{\bf y})}{\partial y_{jt}}}{q_{kmt}({\bf p},{\bf y})}
	\end{eqnarray*}



\section{Demand Specification Sensitivity}
\label{app:demand-spec}

The main text reports a single demand specification, the nested logit whose estimates appear there. Here we report how those estimates change as we add terms to the utility function and how they respond to the premium control function. Table~\ref{tab:app-demand-spec} reports four specifications, each estimated on the same 20\% sample and the same 114 rating-region-by-year cells as the main-text model. Column (1) is a nested logit with plan attributes and a single premium coefficient. Column (2) adds the demographic heterogeneity in the price coefficient and in the actuarial value coefficient, along with demographic shifters of the value of enrolling. Column (3) adds the assistance and commission steering terms and is the specification used in the main text, reproducing those demand estimates exactly. Column (4) adds the premium control function.

We summarize price sensitivity by the enrollment-weighted mean own-price elasticity rather than by the premium coefficient. The premium coefficient and the nesting parameter are jointly identified and trade off against one another as terms are added, so neither is stable on its own. The nesting parameter rises from 0.13 to 0.37 across the columns and the premium coefficient more than doubles in magnitude, but the own-price elasticity they jointly imply stays between $-1.32$ and $-1.13$. The elasticity is measured with respect to the net premium households pay.

We likewise report the assistance results as marginal effects on plan choice rather than as log-odds coefficients. Assigning navigator assistance to every household raises the share choosing a silver plan by 3.2 percentage points, and agent assistance raises it by 2.7 percentage points. The commission-by-agent coefficients are 0.027 on the linear term and $-0.066$ on the squared term.

Columns (3) and (4) differ only in the premium control function. Adding it leaves the nesting parameter, both assistance effects, and the commission coefficients nearly unchanged, and moves the own-price elasticity from $-1.13$ to $-1.25$. \citet{saltzman2019} similarly finds that this correction leaves the nesting parameter essentially unchanged, with only a small downward bias in premium sensitivity when it is omitted.

\begin{table}[!h]
\centering
\begin{minipage}[h]{6in}
\caption[]{\textbf{Demand Specification Sensitivity}\footnote{Four nested-logit demand specifications, each estimated on the 20\% sample and the 114 rating-region-by-year cells used in the main text. Column (1) is a nested logit with plan attributes and a single premium coefficient; column (2) adds demographic heterogeneity in the price coefficient and in the valuation of actuarial value, and demographic shifters of the value of enrolling; column (3) adds the assistance and commission steering terms and reproduces the main-text demand estimates; column (4) adds the premium control function. The mean own-price elasticity is the enrollment-weighted average with respect to the net premium; the navigator and agent effects are the average marginal effect of assistance on the probability of choosing a silver plan, in percentage points. Empty cells denote terms not included in that specification.}}
\label{tab:app-demand-spec}
\centerline{\begin{tabular}{lcccc}
\hline\hline
 & (1) & (2) & (3) & (4) \\
\hline
Mean own-price elasticity & -1.32 & -1.18 & -1.13 & -1.25 \\
$\lambda$ (nesting parameter) & 0.13 & 0.37 & 0.37 & 0.35 \\
Navigator effect on silver (pp) &  &  & 3.2 & 3.0 \\
Agent effect on silver (pp) &  &  & 2.7 & 2.7 \\
Commission $\times$ agent &  &  & 0.027 & 0.026 \\
Commission$^2$/100 $\times$ agent &  &  & -0.066 & -0.070 \\
\hline\hline
\end{tabular}
}
\end{minipage}
\end{table}

\section{Inertia}
\label{app:inertia}

Inertia is a well-documented feature of health insurance plan choice, in employer coverage \citep{handel2013}, in Medicare Part D \citep{polyakova2016}, and in the ACA exchanges we study \citep{saltzman2025}; however, as discussed in the main text, it is not clear how inertia should be treated in a model of commissions and potential steering. For example, a household steered into a plan in year $t$ will tend to stay in that plan in year $t+1$, so that the effect of commissions persists across future years. Including an indicator for prior-year plan, as is often used to capture inertia, would therefore remove the persistence of any steering effects.

Nonetheless, the exclusion of inertia does not appear to meaningfully affect our main estimates. First, in the design-based analysis, the effects of assistance on dominated-plan choice for new enrollees alone (Table~\ref{tab:app-dominated-sample}), for whom inertia is irrelevant, are at least as large as in the full sample for every channel. Second, Table~\ref{tab:app-inertia-lag} re-estimates the primary demand specification with an indicator for the household's prior-year plan, which enters plan utility and therefore the enrollment decision. We estimate a large switching cost of about \$4,200 per enrollee per year; however, the commission coefficients scaled by the nesting parameter are largely unchanged from the main-text estimates. Given the dynamic complexity of inertia in a model of potential steering, and the insensitivity of our main estimates to the inclusion or exclusion of lagged plan indicators, we do not incorporate inertia separately into our estimates.

\begin{table}[!h]
\centering
\begin{minipage}[h]{6in}
\caption[]{\textbf{Demand Estimates with the Prior-Year Plan}\footnote{Main-text demand specification estimated on the 20\% structural sample, with and without an indicator for the household's plan in the prior year. The indicator is zero in 2014 and for households without a prior-year exchange plan. Premiums are net of subsidy and measured in hundreds of dollars.}}
\label{tab:app-inertia-lag}
\centerline{\begin{tabular}{lcc}
\hline\hline
 & Body & With prior-year plan \\
\hline
Premium & -0.4146 & -0.6177 \\
Actuarial value (AV) & 4.3917 & 7.1537 \\
Navigator $\times$ AV & 1.1850 & 2.9859 \\
Agent $\times$ AV & 0.7291 & 1.3968 \\
Navigator $\times$ premium & -0.0976 & -0.3102 \\
Agent $\times$ premium & -0.0214 & -0.1012 \\
Commission $\times$ agent & 0.0273 & 0.0514 \\
Commission$^2$/100 $\times$ agent & -0.0658 & -0.1071 \\
Prior-year plan &  & 3.4395 \\
$\lambda$ (nesting parameter) & 0.3736 & 0.7107 \\
\hline\hline
\end{tabular}
}
\end{minipage}
\end{table}

\section{Cost Sharing and the Objective Welfare Measure}
\label{app:welfare}

The objective welfare measure in the main text values each insured plan at the negative of the household's annual premium plus the mean-variance certainty equivalent of its out-of-pocket exposure. We state here the cost-sharing schedule and the out-of-pocket calculation behind that measure.

Because Covered California standardizes its benefit designs, a plan's cost-sharing is fixed by its metal tier and cost-sharing-reduction level rather than estimated. Table~\ref{tab:app-cost-sharing} reports the 2019 schedule. The full 2014--2019 schedule, which the welfare computation uses, follows the same tier structure, with deductibles and out-of-pocket maxima that rose over the period.

\begin{table}[!h]
\centering
\begin{minipage}[h]{6in}
\caption[]{\textbf{Covered California Standardized Cost Sharing, 2019}\footnote{Standardized patient-centered benefit designs for the 2019 plan year. AV is the statutory actuarial value; the deductible and maximum out-of-pocket (MOOP) are annual individual amounts; coinsurance is the enrollee share of spending above the deductible. CSR denotes the cost-sharing-reduction silver variants available to enrollees below 250\% of the federal poverty level, and the bronze health-savings-account variant is shown separately. The full 2014--2019 schedule follows the same tier structure and is used in the welfare computation.}}
\label{tab:app-cost-sharing}
\centerline{\begin{tabular}{lrrrr}
\hline\hline
Tier & AV & Deductible & Coins. & MOOP \\
\hline
Bronze & 0.60 & \$6,300 & 100\% & \$7,550 \\
Bronze (HSA) & 0.60 & \$6,000 & 40\% & \$6,650 \\
Silver & 0.70 & \$2,500 & 20\% & \$7,550 \\
Silver (CSR 73) & 0.73 & \$2,200 & 20\% & \$6,300 \\
Silver (CSR 87) & 0.87 & \$650 & 15\% & \$2,600 \\
Silver (CSR 94) & 0.94 & \$75 & 10\% & \$1,000 \\
Gold & 0.80 & \$0 & 20\% & \$7,200 \\
Platinum & 0.90 & \$0 & 10\% & \$3,350 \\
\hline\hline
\end{tabular}
}
\end{minipage}
\end{table}

For a plan with deductible $d$, coinsurance rate $c$, and maximum out-of-pocket $M$, out-of-pocket spending at realized annual medical spending $s$ is
\begin{equation}
\label{eqn:app-oop}
o(s) = \min\!\left\{\min(s,\,d) + c\,\max(s-d,\,0),\; M\right\}.
\end{equation}
The household pays all spending up to the deductible, the coinsurance share of spending above it, and never more than the maximum out-of-pocket.

We treat annual individual spending as lognormal with mean $\bar{s}$ and coefficient of variation $\kappa$. Writing $\sigma^2 = \log(1+\kappa^2)$, the underlying normal has variance $\sigma^2$ and mean $\log \bar{s} - \sigma^2/2$, so the spending distribution has mean $\bar{s}$ by construction. We integrate the schedule in Equation~\eqref{eqn:app-oop} against this distribution by deterministic quadrature over 400 evenly spaced midpoint quantiles, which fixes the mean and variance of out-of-pocket spending, $\mathbb{E}[o]$ and $\mathrm{Var}[o]$, without simulation noise. The plan's objective value is then
\begin{equation}
\label{eqn:app-vobj}
V^{obj} = -\left(P + \mathbb{E}[o] + \tfrac{\rho}{2}\,\mathrm{Var}[o]\right),
\end{equation}
with $P$ the annual premium and $\rho$ the coefficient of absolute risk aversion. The variance term gives cost-sharing depth weight beyond its effect on expected out-of-pocket spending, so a household will pay more than the actuarial difference to move from a bronze plan to a platinum one.

The mean $\bar{s}$ is household-specific, set from Medical Expenditure Panel Survey total-expenditure cell means by age composition and income bracket, so that expected spending follows the household rather than the plan and the measure carries no moral-hazard response. The coefficient of variation $\kappa$ and the risk-aversion coefficient $\rho$ are the two calibrated inputs, set to 2.5 and $2.3\times10^{-4}$ per dollar as described in the main text, the latter the marketplace benchmark in \citet{handel2013}. The uninsured option is valued separately, since it carries no plan cost-sharing, using the realized out-of-pocket and social-cost construction described in the main text.

\clearpage

\bibliographystyle{aer}
\bibliography{BibTeX_Library}

\end{document}
